% Section 6.8, from lectures/L06/README.md.

\section{Summary}
\label{c6:sec:review}

\needspace{6\baselineskip}
\subsection*{What you should be able to do now}
\begin{itemize}
  \item Build a 4-bit counter as a gate network, and explain, with your finger on your own drawing,
    why nothing in it makes the count return to zero.
  \item Implement the same counter in VHDL, explain the wraparound as a consequence of the
    register's bit width, and say why a simulator objects to a wraparound that the hardware performs
    without blinking.
  \item Implement a shift register in both SIPO and PISO form, know which suits which job, and know
    which direction the book shifts in and why sticking to one matters.
  \item Combine a counter with a shift register to detect that a complete word has arrived, and
    verify a clocked design by running its testbench.
\end{itemize}

\needspace{6\baselineskip}
\subsection*{Questions to test yourself with}
\begin{enumerate}
  \item Why does a counter of type \code{natural range 0 to 15} return to \code{0} after \code{15}
    with no explicit ``if at the maximum, clear'' logic? What has to hold for the range? Which
    element performs the wraparound?
  \item A VHDL simulator aborts on the wraparound your CircuitVerse drawing performs without
    complaint. Which of the two models the hardware, and what does that say about trusting either
    one alone?
  \item What is the practical difference between a SIPO and a PISO shift register? Which would you
    reach for to receive a stream of serial bits, and which to drive a chain of LEDs from a fixed
    pattern?
  \item A shift register and a counter are both ``N flip-flops sharing a clock''. What separates
    them is what feeds each flip-flop's input. Describe that difference for each.
  \item Why is \code{data_ready} in \code{serial_rx8} a one-clock-cycle pulse rather than a level
    that stays high once a byte has arrived?
  \item Why must \code{shift_enable} gate the bit counter and not just the shifting?
  \item The board demonstration drives \code{shift_enable} from a button instead of slowing the
    clock down. Why is a slowed clock the wrong answer, and what would break if the button were
    connected to \code{shift_enable} without passing through \code{button_sync} first?
  \item \code{data_ready} needs a flip-flop of its own before an LED can show it. Which chapter's
    problem is that, in another guise?
\end{enumerate}

\needspace{6\baselineskip}
\subsection*{Looking ahead}
\Chapref{c7:ch} goes on with:
\begin{itemize}
  \item Timers: a counter with a target value, which is what turns clock cycles into real time.
  \item Built by adding exactly two things to the counter you drew here, a comparator and a clear,
    so keep that CircuitVerse project.
  \item The walking LED: this chapter's shift register, clocked by that timer and started by a
    synchronised button.
\end{itemize}
